%\documentclass[a4paper,10pt,oneside]{article}
%\usepackage{xeCJK}
%\setCJKmainfont{微軟正黑體}
%\begin{document}

\subsubsection{Pick公式}
給定頂點坐標均是整點的簡單多邊形，面積=內部格點數+邊上格點數/2-1

\subsubsection{圖論}
\begin{enumerate}\itemsep = -5pt
\item 對於平面圖，$F=E-V+C+1$，C是連通分量數
\item 對於平面圖，$E\leq 3V-6$
\item 對於連通圖G，最大獨立點集的大小設為I(G)，最大匹配大小設為M(G)，最小點覆蓋設為Cv(G)，最小邊覆蓋設為Ce(G)。
對於任意連通圖：
	\begin{enumerate}\itemsep = -3pt
	\item $I(G)+Cv(G)=|V|$
	\item $M(G)+Ce(G)=|V|$
	\end{enumerate}
\item 對於連通二分圖：
	\begin{enumerate}\itemsep = -3pt
	\item $I(G)=Cv(G)$
	\item $M(G)=Ce(G)$
	\end{enumerate}
\end{enumerate}

\subsubsection{學長公式}
\begin{enumerate}\itemsep = -3pt
\item $\sum_{d|n} \phi(n) = n$
\item $\text{Harmonic series } H_n = \ln(n) + \gamma + 1/(2n) - 1/(12n^2) + 1/(120n^4)$
\item $ \gamma =$ 0.57721566490153286060651209008240243104215
\item 選轉矩陣$M(\theta)= \left( \begin{array}{ccc}
cos\theta & -sin\theta \\ 
sin\theta &  cos\theta
\end{array} \right)$
\end{enumerate}

\subsubsection{基本數論}
\begin{enumerate}\itemsep = -3pt
	\item $\sum_{d|n} \mu (n)=[n==1]$
	\item $g(m)=\sum_{d|m}f(d)\Leftrightarrow f(m)=\sum_{d|m}\mu (d) \times g(m/d)$
	\item $\sum_{i=1}^n\sum_{j=1}^m$互質數量$=\sum \mu (d)\left \lfloor \frac{n}{d} \right \rfloor \left \lfloor \frac{m}{d} \right \rfloor$
	\item $\sum_{i=1}^n\sum_{j=1}^nlcm(i,j)=n\sum_{d|n} d \times \phi (d)$
\end{enumerate}

\subsubsection{排組公式}
\begin{enumerate}
\itemsep = -3pt
\item k卡特蘭$\frac{C_{n}^{kn}}{n(k-1)+1}$，$C_{m}^{n}=\frac{n!}{m!(n-m)!}$
\item $H(n,m)\cong x_1+x_2\ldots +x_n=k, num=C_{k}^{n+k-1}$
\item Stirling number of $2^{nd}$,$n$人分$k$組方法數目
	\begin{enumerate}\itemsep = -2pt
		\item $S(0,0)=S(n,n)=1$
		\item $S(n,0)=0$
		\item $S(n,k)=kS(n-1,k)+S(n-1,k-1)$
	\end{enumerate}
\item Bell number,$n$人分任意多組方法數目
	\begin{enumerate}\itemsep = -2pt
		\item $B_0=1$
		\item $B_n=\sum_{i=0}^nS(n,i)$
		\item $B_{n+1}=\sum_{k=0}^{n} C_k^n B_k$
		\item $B_{p+n}\equiv B{_n}+B_{n+1} mod p$, p is prime
		\item $B_{p^m+n}\equiv mB{_n}+B_{n+1} mod p$, p is prime
		\item From $B_0: 1,1,2,5,15,52,\\203,877,4140,21147,115975$
	\end{enumerate}
\item Derangement,錯排,沒有人在自己位置上
	\begin{enumerate}\itemsep = -2pt
		\item $D_n=n!(1-\frac{1}{1!}+\frac{1}{2!}-\frac{1}{3!}\ldots +(-1)^n\frac{1}{n!})$
		\item $D_n=(n-1)(D_{n-1}+D_{n-2}),D_0=1,D_1=0$
		\item From $D_0: 1,0,1,2,9,44,\\265,1854,14833,133496$
	\end{enumerate}
\item Binomial\ Equality
	\begin{enumerate}\itemsep = -2pt
	    \item $\sum_k \binom{r}{m + k} \binom{s}{n - k} = \binom{r + s}{m + n}$
         \item $\sum_k \binom{l}{m + k} \binom{s}{n + k} = \binom{l + s}{l -m + n}$		
         \item $\sum_k \binom{l}{m + k} \binom{s + k}{n}(-1)^k = (-1)^{l + m} \binom{s - m}{n - l}$
		\item $\sum_{k\leq l} \binom{l - k}{m} \binom{s}{k - n}(-1)^k = (-1)^{l + m} \binom{s - m - 1}{l - n - m}$
		\item $\sum_{0 \leq k \leq l} \binom{l - k}{m} \binom{q + k}{n} = \binom{l + q + 1}{m + n + 1}$
		\item $\binom{r}{k} = (-1)^k\binom{k - r - 1}{k}$
		\item $\binom{r}{m} \binom{m}{k} = \binom{r}{k} \binom{r - k}{m - k}$
		\item $\sum_{k\leq n} \binom{r + k}{k} = \binom{r + n + 1}{n}$
		\item $\sum_{0\leq k \leq n} \binom{k}{m} = \binom{n + 1}{m + 1}$
		\item $\sum_{k\leq m}\binom{m + r}{k}x^ky^k = \sum_{k\leq m}\binom{-r}{k}(-x)^k (x+y)^{m-k}$	
	\end{enumerate}
\end{enumerate}

\subsubsection{冪次,冪次和}
\begin{enumerate}\itemsep = -3pt
	\item $a^b\%P=a^{b\% \varphi (p)+\varphi (p)},b\geq \varphi (p)$
	\item $1^3+2^3+3^3+\ldots +n^3=\frac{n^4}{4}+\frac{n^3}{2}+\frac{n^2}{4}$
	\item $1^4+2^4+3^4+\ldots +n^4=\frac{n^5}{5}+\frac{n^4}{2}+\frac{n^3}{3}-\frac{n}{30}$
	\item $1^5+2^5+3^5+\ldots +n^5=\frac{n^6}{6}+\frac{n^5}{2}+\frac{5n^4}{12}-\frac{n^2}{12}$
	\item $0^k+1^k+2^k+\ldots +n^k = P(k),P(k)=\frac{(n+1)^{k+1}-\sum_{i=0}^{k-1}C_i^{k+1}P(i)}{k+1},P(0)=n+1$
	\item $\sum_{k=0}^{m-1}k^n=\frac{1}{n+1}\sum_{k=0}^{n}C_k^{n+1}B_km^{n+1-k}$
	\item $\sum_{j=0}^{m}C_j^{m+1}B_j=0,B_0=1$
	\item 除了$B_1=-1/2$，剩下的奇數項都是$0$
	\item $B_2=1/6,B_4=-1/30,B_6=1/42,B_8=-1/30,B_{10}=5/66,B_{12}=-691/2730,B_{14}=7/6,B_{16}=-3617/510,B_{18}=43867/798,B_{20}=-174611/330,$
\end{enumerate}

\subsubsection{循環小數轉分數}
\begin{enumerate}\itemsep = -3pt
\item 若$x=0.\overline{a}$，則
\[
x = \frac{a}{\underbrace{99\ldots9}_{k \text{ digits}}}
\]
其中 $a$ 為循環節，$k$ 為循環節的位數。
\item 例子：
\[
0.\overline{37} = \frac{37}{99}
\]
\[
0.\overline{5} = \frac{5}{9}
\]
\end{enumerate}

\subsubsection{循環小數轉分數}
\begin{enumerate}\itemsep = -3pt
\item 純循環小數：若 $x = 0.\overline{a}$，其中 $a$ 為循環節、長度為 $k$，
\[
x = \frac{a}{\underbrace{99\ldots9}_{k \text{ digits}}}
\]
例：
\[
0.\overline{37} = \frac{37}{99}, \quad 0.\overline{5} = \frac{5}{9}
\]

\item 混循環小數：若 $x = 0.b\overline{a}$，其中 $b$ 為前綴、長度 $m$，$a$ 為循環節、長度 $k$，
\[
x = \frac{(b \cdot 10^k + a) - b}{10^m(10^k-1)}
\]

例：
\[
0.12\overline{3} 
= \frac{(12 \cdot 10^1 + 3) - 12}{10^2(10^1-1)}
= \frac{123 - 12}{100 \cdot 9}
= \frac{111}{900}
= \frac{37}{300}
\]
\end{enumerate}

\subsubsection{常見級數與組合公式}
\begin{enumerate}\itemsep = -3pt
\item 平方和公式：
\[
1^2 + 2^2 + \cdots + n^2 = \frac{n(n+1)(2n+1)}{6}
\]
\[
4^2 + 6^2 + \cdots + (2n)^2 = \frac{(2n)(n+1)(2n+1)}{3}
\]
\[
1^2 + 3^2 + \cdots + (2n+1)^2 = \frac{n(2n-1)(2n+1)}{3}
\]

\item 立方和公式：
\[
1^3 + 2^3 + \cdots + n^3 = \frac{n^4 + 2n^3 + n^2}{4}
\]

\item 四次方和公式：
\[
1^4 + 2^4 + \cdots + n^4 = \frac{6n^5 + 15n^4 + 10n^3 - n}{30}
\]

\item 五次方和公式：
\[
1^5 + 2^5 + \cdots + n^5 = \frac{2n^6 + 6n^5 + 5n^4 - n^2}{12}
\]

\item 六次方和公式：
\[
1^6 + 2^6 + \cdots + n^6 = \frac{6n^7 + 21n^6 + 21n^5 - 7n^3 + n}{42}
\]

\item 七次方和公式：
\[
1^7 + 2^7 + \cdots + n^7 = \frac{3n^8 + 12n^7 + 14n^6 - 7n^4 + 2n^2}{24}
\]	

\item 八次方和公式：
\[
\begin{aligned}
1^8 + 2^8 + \cdots + n^8 \\
= \frac{10n^9 + 45n^8 + 60n^7 - 42n^5 + 20n^3 - 3n}{90}
\end{aligned}
\]

\item 九次方和公式：
\[
\begin{aligned}
1^9 + 2^9 + \cdots + n^9 \\
= \frac{2n^{10} + 10n^9 + 15n^8 - 14n^6 + 10n^4 - 3n^2}{20}
\end{aligned}
\]

\item 十次方和公式：
\[
\begin{aligned}
1^{10} + 2^{10} + \cdots + n^{10} \\
= \frac{6n^{11} + 33n^{10} + 55n^9 - 66n^7 + 66n^5 - 33n^3 + 5n}{66}
\end{aligned}
\]

\item 等比級數：
\[
S = a \cdot \frac{r^n - 1}{r - 1}
\]

\item 二項式係數恆等式：
\[
\binom{n}{0} + \binom{n}{1} + \cdots + \binom{n}{n} = 2^n
\]
\[
\binom{n}{1} + \binom{n}{2} + \cdots + \binom{n}{n} = 2^n - 1
\]
\[
\binom{n}{1} + \binom{n}{3} + \cdots + \binom{n}{n - k} = 2^{n - 1}
\]

\item 分配問題（玩具分給小孩）：
\begin{enumerate}\itemsep = -3pt
\item $n$ 個玩具，$k$ 位小孩，可以有人沒拿到：
\[
\binom{n+k-1}{n} = \binom{n+k-1}{k-1}
\]
\item $n$ 個玩具，$k$ 位小孩，每個人至少一個：
\[
\binom{n-1}{k-1}
\]
\end{enumerate}

\subsubsection{位元運算}
\begin{enumerate}\itemsep = -3pt
\item 位元條件：
\[
(x + k) \,\&\, (y + k) = 0
\]

\item 加法恆等式（利用 XOR 與 AND）：
\[
a + b = (a \oplus b) + 2 \cdot (a \,\&\, b)
\]
\item OR 與 AND 的關係：
\[
a \,|\, b = a + b - (a \,\&\, b)
\]
\item 交換兩數：
\[
a = a \oplus b, \quad b = a \oplus b, \quad a = a \oplus b
\]
\item 取得最低位的 $1$：
\[x \,\&\, (-x)
\]
\item 清除最低位的 $1$：
\[x \,\&\, (x - 1)
\]
\end{enumerate}

\subsubsection{除數與倍數}
\begin{enumerate}\itemsep = -3pt
\item LCM:
\[
\text{lcm}(a, b) = \frac{a \times b}{\gcd(a, b)}
\]

\item LCM (prime factorization):
\[
\text{lcm}(a, b) = \prod_{p \in \text{primes}} p^{\max(e_a, e_b)}
\]

\item GCD (prime factorization):
\[
\gcd(a, b) = \prod_{p \in \text{primes}} p^{\min(e_a, e_b)}
\] 

\item Counting divisors:
\[
\text{If } n = \prod_{i=1}^{k} p_i^{e_i},  cnt \text{ is } \prod_{i=1}^{k} (e_i + 1)
\]
\end{enumerate}

\subsubsection{數論公式}
\item Bezout's identity：
\[
ax + by = \gcd(a, b) \quad \text{（必定存在整數解 $x, y$）}
\]

\item 模指數運算（冪的冪）：
\[
a^{\,b^c} \equiv a^{\,\text{power\_mod}(b, c, \, \text{MOD}-1)} \pmod{\text{MOD}}
\]
\end{enumerate}

\item Pythagorean theorem：
\[
(n^2 + m^2, 2nm, n^2 - m^2), \quad n > m > 0
\]

\item Wilson's theorem：
\[
(p-1)! \pmod{p} = n-1 \quad \text{（$p$ 為質數）}
\]

\item Catalan number：
\[
C_n = \frac{1}{n+1} \binom{2n}{
n} = \frac{(2n)!}{(n+1)!n!}
\]

\item 海龍公式：
\[
\text{Area} = \sqrt{s(s-a)(s-b)(s-c)}, \quad s = \frac{a+b+c}{2}
\]

\item 兩圓相交弦長：
\[
\text{Length} = 2 \cdot \sqrt{\frac{(s-a)(s-b)}{s}}, \quad s = \frac{r_1 + r_2 + d}{2}
\]

\item 三角形內切圓半徑：
\[
r = \frac{2 \cdot \text{Area}}{a + b + c}
\]

\item 三角形外接圓半徑：
\[
R = \frac{abc}{4 \cdot \text{Area}}
\]

\item 三角形重心：
\[
G = \left(\frac{x_1 + x_2 + x_3}{3}, \frac{y_1 + y_2 + y_3}{3}\right)
\]

\item 三角形垂心：
\[
H = \left(\frac{x_1 \tan A + x_2 \tan B + x_3 \tan C}{\tan A + \tan B + \tan C}, \frac{y_1 \tan A + y_2 \tan B + y_3 \tan C}{\tan A + \tan B + \tan C}\right)
\]

\item 三角形內心：
\[
I = \left(\frac{a x_1 + b x_2 + c x_3}{a + b + c}, \frac{a y_1 + b y_2 + c y_3}{a + b + c}\right)
\]

\item 奇(偶)長度陣列子數量：
\[
n = \frac{(i + 1) * (n - i) + 1}{2}
\]

%\end{document}